\addtocontents{toc}{\protect\Needspace{3\baselineskip}}
\originalchapter{世界面费米子、GSO 投影与态计数}\label{ch:basic-super}
对应 8.251 第 21 讲后半与第 22 讲. 原课对超弦采用结果与直观解释相结合的方式. 本章补足自由费米场、边界条件、正规序与低层计数; 超共形鬼场、全局模不变性和完整时空超对称代数仍明确属于高阶待补推导.

\originalsection{Grassmann 变量与量子反对易}
经典 Grassmann 奇变量满足 $\theta_i\theta_j=-\theta_j\theta_i$, 因而 $\theta_i^2=0$. 这是定义, 不能用普通实数代替. 量子费米振子满足
\[
 \{f_i,f_j^\dagger\}=\delta_{ij},\qquad
 \{f_i,f_j\}=0=\{f_i^\dagger,f_j^\dagger\}.
\]
于是 $(f_i^\dagger)^2=0$, 每个模式的占据数为 0 或 1. 若只是两块彼此反对易的 Clifford 矩阵, 它们各自平方可为 1, 并不等于经典 Grassmann 变量; 原课的矩阵对照正是用来强调此区别.

对八个横向方向, 引入两个实世界面费米场 $\psi_1^I,\psi_2^I$. 一种方便的无量纲归一化为
\[
 S_f=\frac{\ii}{4\pi}\int\dd\tau\dd\sigma
 [\psi_1^I(\partial_\tau+\partial_\sigma)\psi_1^I
 +\psi_2^I(\partial_\tau-\partial_\sigma)\psi_2^I].
\]
对奇变量变分并将 $\delta\psi$ 放在同一侧, 两次积分分部的反交换符号使一阶项相加. 体方程为 $(\partial_\tau+\partial_\sigma)\psi_1=0$ 与 $(\partial_\tau-\partial_\sigma)\psi_2=0$, 故 $\psi_1=\Psi_1(\tau-\sigma)$、$\psi_2=\Psi_2(\tau+\sigma)$. 空间边界项正比 $[\psi_1\delta\psi_1-\psi_2\delta\psi_2]_0^\pi$.

\originalsection{开弦的两种费米边界条件}
可在两端分别取 $\psi_1=\eta\psi_2$, $\eta=\pm1$, 使上述边界项消失. 一个端点的整体符号可选为正, 两端的相对符号仍有两种不同选择. 将两支拼成 $-\pi\leq\sigma\leq\pi$ 的一个费米場, 两端同号得到周期 Ramond（R）条件, 异号得到反周期 Neveu--Schwarz（NS）条件.

因此模式指数在 R 扇区为整数, 在 NS 扇区为半整数:
\[
 \Psi_{\rm R}=\sum_{n\in\mathbb Z}d_n^I\ee^{-\ii n(\tau-\sigma)},
 \quad\Psi_{\rm NS}=\sum_{r\in\mathbb Z+1/2}b_r^I\ee^{-\ii r(\tau-\sigma)}.
\]
相应正则反对易式为 $\{d_m^I,d_n^J\}=\delta^{IJ}\delta_{m+n,0}$、
$\{b_r^I,b_s^J\}=\delta^{IJ}\delta_{r+s,0}$. 正频率为湮灭算符, 负频率为创建算符. 不应在两端分别令一个左行费米场恒为零: 第一阶传播方程会把这种条件扩展为整支场消失.

\originalsection{费米零点能及十维自由质量式}
每个玻色方向的有限零点能为 $-1/24$. 半整数正频率的调节有限和为
\[
 \sum_{r>0, r\in\mathbb Z+1/2}r\Big|_{\rm finite}
 =\zeta(-1,\tfrac12)=\frac1{24}.
\]
费米子因反对易正规序带负号, 乘 $-1/2$ 得每个 NS 方向 $-1/48$. 八个方向合起来 $8(-1/24-1/48)=-1/2$. R 正整数费米零点能为 $+1/24$, 与玻色逐方向抵消. 因而
\[
 \alpha'M_{\rm NS}^2=N_B+N_F-\tfrac12,
 \qquad\alpha'M_{\rm R}^2=N_B+N_F,
\]
$N_F=\sum_{r>0}r b_{-r}^Ib_r^I$ 或 $\sum_{n>0}n d_{-n}^Id_n^I$. R 零模不贡献此层数, 却决定真空简并和自旋.

十维临界条件还可在协变 RNS 体系中表示为物质中央荷 $c_{\rm matter}=D+D/2$ 与两类规范鬼场的总 $c_{\rm ghost}=-15$ 抵消, 因而 $3D/2=15$. 这里列出的是已建立的鬼场中央荷结果, 并未在本入门章重建 BRST 与超共形测度. 仅在任意 $D$ 里写同样振子求和不保证目标洛伦兹一致性; 第 26 讲的 $D=18$ 练习应理解为形式计数检查.

\originalsection{Ramond 零模为何给自旋态}
R 零模有 $\{d_0^I,d_0^J\}=\delta^{IJ}$. 令 $\gamma^I=\sqrt2d_0^I$, 得 $SO(8)$ 的 Clifford 代数 $\{\gamma^I,\gamma^J\}=2\delta^{IJ}$. 将八个实零模组合成四对创建、湮灭算符, 基态空间维数为 $2^4=16$. 以零模费米数的奇偶分开, 偶数子空间有 $1+6+1=8$ 态, 奇数有 $4+4=8$ 态, 分别为两个手征旋量 $8_s,8_c$.

零模真空的时空自旋是旋量, 不是向量. 因而 R 扇区最终提供时空费米子, 尽管原来的 $\psi^I$ 在世界面上携带目标横向指标. 世界面费米性与时空玻色、费米统计不能单靠名字判断, 要看量子零模的小群表示.

\originalsection{GSO 投影与最低层}
在本章的振子费米数约定下, NS 保留奇数个费米创建算符的态, 去掉 $N=0$ 的快子, 留下 $b_{-1/2}^I|\mathrm{NS}\rangle$ 的八个无质量向量极化. R 选择一个总手征, 留下八个无质量旋量极化. 非零费米激发改变手征奇偶, 因此某 massive 层的允许零模手征须和振子奇偶共同判断.

第一 massive NS 层 $N=3/2$ 有 $b_{-3/2}^I$ 的 8 态、
$\alpha_{-1}^Ib_{-1/2}^J$ 的 64 态、$b_{-1/2}^{[I}b_{-1/2}^Jb_{-1/2}^{K]}$ 的 56 态, 总数 128. 第一 R massive 层则有 $\alpha_{-1}^I|s\rangle$ 与 $d_{-1}^I|s'\rangle$, 各 64 态, 同样 128. 这比仅说“超对称保证相等”更具体, 但只检验了此层, 完整对称性还需一般代数或生成函数恒等式.

\originalsection{生成函数将重复计数变成乘法}
一个能级贡献 $n$ 的玻色模式可重复激发, 因子为 $(1-q^n)^{-1}$; 一个费米模式最多激发一次, 因子为 $1+q^n$. 因而临界开玻色弦的单弦内部态计数是
\[
 Z_B(q)=q^{-1}\prod_{n=1}^\infty(1-q^n)^{-24}.
\]
前面的 $q^{-1}$ 使指数直接等于 $\alpha'M^2$, 去掉该项时指数才是层数 $N$. 例如 $N=4$ 的整数分拆为 $4,3+1,2+2,2+1+1,1+1+1+1$, 不同模式可独立赋横向指标, 同一模式则只数对称乘积.

GSO 后的十维开超弦生成函数为
\begin{align*}
 Z_{\rm NS}(q)&=\frac1{2\sqrt q}\prod_{n=1}^\infty(1-q^n)^{-8}
 \left[\prod_{n=1}^\infty(1+q^{n-1/2})^8
 -\prod_{n=1}^\infty(1-q^{n-1/2})^8\right],\\
 Z_{\rm R}(q)&=8\prod_{n=1}^\infty\left(\frac{1+q^n}{1-q^n}\right)^8.
\end{align*}
两者相等是 Jacobi 恒等式的等价形式, 本章引用此一般恒等式. 随包程序独立枚举有限低层核验, 前五项依次为 8、128、1152、7680、42112. 这些是每个扇区的数目, 同时把 NS 与 R 合起来会再乘二.

\originalsection{指数增长与 Hagedorn 温度}
对 $k$ 个玻色方向, 令 $q=\ee^{-\beta}$ 且 $\beta\ll1$, 有
\[
 \log\prod_n(1-q^n)^{-k}
 =k\sum_{m,n\geq1}\frac{\ee^{-\beta mn}}m
 \simeq\frac{k\pi^2}{6\beta}.
\]
Laplace 反演的指数 $\beta N+k\pi^2/(6\beta)$ 在
$\beta_* =\pi\sqrt{k/(6N)}$ 驻定, 给 $\log d_N\sim2\pi\sqrt{kN/6}$. 一维分拆的更精细前因子不等于全体 24 方向的前因子, 不应照搬. 对开玻色弦 $k=24$, 高质量 $M\sim\sqrt{N}/\ell_s$, 状态密度按 $\exp(4\pi\ell_sM)$ 增长, 自由正则配分函数的收敛温度尺度是 $T_H=1/(4\pi\ell_s)$.

对超弦, 八个玻色与八个费米模式的有效中央荷为 12, 相同估算给 $T_H=1/(2\sqrt2\pi\ell_s)$. 这是自由弦的谱增长和热配分函数收敛条件, 不能单独判断相互作用后的实际热相变或高温终态.
